\documentclass[11pt]{article}
\usepackage[margin=2.3cm]{geometry}

\title{Prueba del visor de salidas (tabla + gráfico)}
\author{}
\date{\today}

\begin{document}
\maketitle

Ejecuta las celdas con \textbf{Cálculo → Todas las celdas} (o
\textbf{Compilar}). En la salida de cada celda, pasa el ratón por encima de la
tabla o del gráfico y pulsa el botón de \emph{ampliar} (esquina superior
derecha) para abrirlo en el visor:

\begin{itemize}
  \item \textbf{Gráfico:} se abre con zoom (rueda), arrastre para desplazar y
        lectura de coordenadas en píxeles bajo el cursor.
  \item \textbf{Tabla:} se abre a tamaño completo, con scroll y texto
        seleccionable para verificar los datos.
\end{itemize}

\section{Tabla generada con pandas}
La salida se muestra como tabla HTML, igual que en Jupyter:

%#python
import numpy as np
import pandas as pd

n = np.arange(1, 16)
df = pd.DataFrame({
    "n": n,
    "n^2": n**2,
    "raiz(n)": np.sqrt(n).round(4),
    "ln(n)": np.log(n).round(4),
    "1/n": (1 / n).round(4),
})
df
%#end

\section{Gráfico con matplotlib}
La figura abierta se captura automáticamente y se muestra en la celda:

%#python
import numpy as np
import matplotlib.pyplot as plt

x = np.linspace(0, 2 * np.pi, 400)
plt.figure(figsize=(6.5, 3.4))
plt.plot(x, np.sin(x), label="sin x")
plt.plot(x, np.cos(x), label="cos x")
plt.plot(x, np.sin(x) * np.cos(x), "--", label="sin x · cos x")
plt.title("Funciones trigonométricas")
plt.xlabel("x"); plt.ylabel("y")
plt.legend(); plt.grid(True, alpha=0.4)
%#end

\medskip
Todo se ejecuta en un kernel de Python real, así que funciona cualquier
librería instalada (puedes añadir más desde \textbf{Ver → Terminal}).

\end{document}
